\scanpage{147}
% archive foliation: 146
% Pale reversed typescript is show-through from the reverse, not writing on this face.
\deleted{Si on a choisi les $\overline{s}''$ et $\overline{s}'$ tels que
\unsure{diag}}
\deleted{\inserted{$s''=\operatorname{diag}_{\unsure{\theta_2}}(\overline{s}')
\text{ et }\deleted{\unsure{s''}}\
s'''=\operatorname{diag}_{\unsure{\theta_3}}(\overline{s}')$}}

\unsure{Rem.} Il est loisible de supposer que pour $\forall\ s'\in E'$,
désignant par $s''$ et $s'''$ ses images dans $E''$, $\unsure{E'''}$
par l'application diagonale,
\deleted{\unsure{on a $\overline{s}''=\operatorname{diag}$}},
$\unsure{\overline{s}'''}$ sont les images de $\overline{s}'$ et
$\overline{s}''$ par les morphismes diagonaux, et qu'on ait choisi
les \unsure{derniers} chemins de \unsure{jonction}, associés à ces $s''$, $s'''$
\deleted{\unsure{comme}} les chemins triviaux.
\inserted{\unsure{Enfin, on peut supposer que $s''=\overline{s}''$.}}
Alors il résulte de c qu'on a
\[
g_{s''}=g_{s''}g_{s''}\qquad\text{i.e.}\qquad g_{s''}=1
\]
et la relation b) prend la forme équivalente
\[
\text{b')}\qquad g_{\overline{s}'}=1
\]

\figure{\begin{tikzpicture}
\draw (0,0) -- (8,0.5);
\draw[line width=2pt] (7.8,-0.1) .. controls (8.6,0.2) and (9.2,0.9) .. (9.4,1.4);
\end{tikzpicture}}

\diagram{\begin{tikzcd}
X' & G.\pi_1 \arrow[dr,bend right=25]\\
&& \phantom{S}
\end{tikzcd}}
\[
S\longleftarrow S'
\]
\[
\pi\qquad\qquad S
\]
\[
\pi_1(S,a)\quad\pi
\]
\diagram{\begin{tikzcd}
 & R_0 \arrow[d] & R_1 \arrow[d] &
R_2 \arrow[d,shift left=4pt] \arrow[d] \arrow[d,shift right=4pt]\\
S & S' \arrow[l] &
S'' \arrow[l,shift left=3pt] \arrow[l,shift right=3pt] &
S''' \arrow[l,shift left=5pt] \arrow[l] \arrow[l,shift right=5pt]
\end{tikzcd}}
